\newcommand{\tipoRequerimento}{Defesa de TOI}

\section{DOS FATOS E DOS DIREITOS}

\begin{enumerate}[resume]
\item Trata-se de uma reclamação referente à unidade consumidora nº \unidadeConsumidora{}.
\item Foi apurado que a Concessionária efetuou uma cobrança, aparentemente relacionada ao consumo não registrado (CNR), resultando na fatura nº 2000033195001 %\fatura{}.
\item Contudo, para que a Concessionária tenha legitimidade para proceder com a referida cobrança, é imprescindível que sejam cumpridas as obrigações regulatórias e legais pertinentes.
\item Ademais, para que o consumidor tenha pleno conhecimento dos valores cobrados e possa exercer adequadamente o seu direito de defesa, é imperioso que a Concessionária remeta os documentos e comprovações estabelecidos no artigo 325, conforme as exigências legais:

\Citacao{
    Art. 325. A distribuidora deve compensar o faturamento quando houver diferença a cobrar ou a devolver decorrente das seguintes situações:

    (..)
    
    II - Comprovação de procedimentos irregulares, de que trata o Capítulo VII do Título II  
    
    (...)
    
    § 1º A distribuidora deve notificar o consumidor por escrito, por modalidade que permita a comprovação do recebimento, contendo obrigatoriamente
    
    (... )
    
    II - No caso de procedimentos irregulares, os itens do caput do art. 598;
}

\item No que tange à documentação relativa à lavratura do TOI, cumpre destacar que é dever da distribuidora manter arquivados todos os documentos pertinentes, incluindo aqueles que deram origem ao referido termo, bem como os demais registros relacionados, conforme estabelece o artigo 598 da Resolução Normativa ANEEL nº 1.000/2021:

\Citacao{
    Art. 598. Nos casos em que houver necessidade de compensação de receita em decorrência da irregularidade apurada, a distribuidora deve instruir um processo com as seguintes informações:

    (..)

    § 1º A distribuidora deve armazenar no processo todas as notificações, reclamações, respostas e outras interações realizadas, bem como demais informações e documentos relacionados ao caso. 

    § 2º O faturamento da compensação deve ser realizado conforme art. 325. (Grifou-se)
}

\item Na hipótese de ser constatada qualquer ocorrência que indique a prática de procedimento irregular, a distribuidora deve emitir o Termo de Ocorrência e Inspeção (TOI), em estrita observância ao disposto no artigo 590 da Resolução Normativa nº 1.000/2021 da ANEEL:

\Citacao{
    Art. 590. Na ocorrência de indício de procedimento irregular, a distribuidora deve adotar as providências necessárias para sua fiel caracterização, compondo um conjunto de evidências por meio dos seguintes procedimentos:
    
    I - Emitir o Termo de Ocorrência e Inspeção – TOI, em formulário próprio, elaborado conforme instruções da ANEEL.
}

\item Verificada a emissão do Termo de Ocorrência e Inspeção (TOI), impõe-se à Concessionária o dever de observar integralmente as disposições contidas no artigo 598 da Resolução Normativa nº 1.000/2021 da ANEEL:

\Citacao{
    Art. 598. Nos casos em que houver necessidade de compensação de receita em decorrência da irregularidade apurada, a distribuidora deve instruir um processo com as seguintes informações:

    I - Ocorrência constatada;

    II - Cópia legível do TOI;

    III - avaliação do histórico de consumo e das demais grandezas elétricas;

    IV - Cópia de todos os elementos de apuração da ocorrência, incluindo as informações da medição fiscalizadora, quando for o caso;

    V - Relatório de avaliação técnica, quando constatada a violação do medidor ou demais equipamentos de medição;

    VI - Comprovantes de notificação, agendamento e reagendamento da avaliação técnica;

    VII - relatório da perícia metrológica, quando solicitada, informando quem solicitou e onde foi realizada;

    VIII - custos de frete e da perícia metrológica, quando esta tiver sido solicitada pelo consumidor e for comprovada a irregularidade;

    IX - Comprovação de que o defeito na medição foi decorrente de aumento de carga à revelia, quando alegado este motivo;

    X - Critério utilizado para a recuperação de receita, conforme art. 595, e a memória descritiva do cálculo realizado, de modo que permita a sua reprodução, e as justificativas para não utilização de critérios anteriores;

    XI - valor do custo administrativo cobrado e o motivo, conforme art. 597;

    XII - critério utilizado para a determinação do período de duração, conforme art. 596, e a memória descritiva da avaliação realizada, de modo que permita a sua reprodução e, quando for o caso, as justificativas pela não adoção dos demais critérios dispostos no artigo;

    XIII - data da última inspeção que antecedeu a inspeção que originou a notificação;

    XIV - valor da diferença a cobrar ou a devolver, com a memória descritiva de como o valor foi apurado; e 

    XV - Tarifas utilizadas.

    [...]

    § 4º A distribuidora deve fornecer em até 5 dias úteis, mediante solicitação do consumidor, a cópia do processo de irregularidade.

    (grifo nosso).
}

\item Nos casos em que a concessionária tenha realizado a lavratura do TOI em decorrência ao defeito na medição é obrigatório que a Concessionária siga os procedimentos expostos no art. 257:

\Citacao{
    Art. 257. Para compensação no faturamento no caso de defeito na medição, a distribuidora deve instruir um processo com as seguintes informações:

    I - Ocorrência constatada;
    
    II - Cópia legível do TOI;
    
    III - os números dos equipamentos e as informações das leituras do medidor retirado e instalado;
    
    IV - Avaliação do histórico de consumo e das demais grandezas elétricas;
    
    V - Relatório da inspeção do sistema de medição, informando as variações verificadas, os limites admissíveis e a conclusão final;
    
    VI - Comprovantes de notificação, agendamento e reagendamento da inspeção;
    
    VII - relatório da verificação do medidor junto ao INMETRO ou órgão delegado, quando solicitada, informando quem solicitou e onde foi realizada;
    
    VIII - custos de frete, da inspeção e verificação atribuíveis ao consumidor e demais usuários;
    
    IX - Critério utilizado para a compensação, conforme art. 255, e a memória descritiva do cálculo realizado, de modo que permita a sua reprodução, e as justificativas para não utilização de critérios anteriores;
    
    X - Critério utilizado para a determinação do período de duração, conforme art. 256;
    
    XI - valor da diferença a cobrar ou a devolver, com a memória descritiva de como o valor foi apurado; e
    
    XII - tarifas utilizadas.
    
    § 2º O faturamento da compensação deve ser realizado conforme art. 325.
    
    § 4º A distribuidora deve fornecer em até 5 dias úteis, mediante solicitação do consumidor, cópia do processo individualizado do defeito na medição.
}

\item Na hipótese da lavratura do Termo de Ocorrência e Inspeção (TOI) ter ocorrido durante a vigência da Resolução Normativa nº 414/2010 da ANEEL, é imperioso que a Concessionária adote a metodologia prevista no artigo 129 do referido diploma normativo, transcrito a seguir:

\Citacao{
    Art. 129. Na ocorrência de indício de procedimento irregular, a distribuidora deve adotar as providências necessárias para sua fiel caracterização e apuração do consumo não faturado ou faturado a menor.

    § 1º A distribuidora deve compor conjunto de evidências para a caracterização de eventual irregularidade por meio dos seguintes procedimentos:
    
    I – Emitir o Termo de Ocorrência e Inspeção – TOI, em formulário próprio, elaborado conforme Anexo V desta Resolução;
    
    II – Solicitar perícia técnica, a seu critério, ou quando requerida pelo consumidor ou por seu representante legal;
    
    III – elaborar relatório de avaliação técnica, quando constatada a violação do medidor ou demais equipamentos de medição, exceto quando for solicitada a perícia técnica de que trata o inciso II; (Redação dada pela REN ANEEL 479, de 03.04.2012)
    
    IV – Efetuar a avaliação do histórico de consumo e grandezas elétricas; e
    
    V – Implementar, quando julgar necessário, os seguintes procedimentos:
    
    a) medição fiscalizadora, com registros de fornecimento em memória de massa de, no mínimo, 15 (quinze) dias consecutivos; e
    
    b) recursos visuais, tais como fotografias e vídeos.
    
    § 2º Uma cópia do TOI deve ser entregue ao consumidor ou àquele que acompanhar a inspeção, no ato da sua emissão, mediante recibo.
}

\item Diante do exposto, verifica-se que a concessionária procedeu à cobrança de forma irregular por ocasião da emissão do Termo de Ocorrência e Inspeção (TOI) referente à Unidade Consumidora acima mencionada, por descumprir todas as determinações acerca da lavratura do TOI conforme estabelecidos no art. 598, e art. 257 da REN 1.000/2021 e art. 129 da REN 414/2010.
\item Caso a Concessionária deixe de apresentar a integralidade dos documentos e informações exigidos no art. 598 e 257 da Resolução Normativa nº 1.000/2021 da ANEEL, bem como no art. 129 da Resolução Normativa nº 414/2010, restarão configurados os fundamentos que ensejam a nulidade do Termo de Ocorrência e Inspeção (TOI) e da respectiva cobrança dele decorrente.
\end{enumerate}

\section{DOS PEDIDOS}

\begin{enumerate}[resume]
\item Ante o exposto, requer-se:
\begin{enumerate}   
\item Solicita-se a análise e o atendimento a esta reclamação no prazo máximo de 5 (cinco) dias úteis, conforme o Art. 408, I, da REN 1000/2021 da ANEEL.
\item Envie cópia do TOI e demais registros e documentos relacionados ao TOI (memorial de cálculo, carta de agendamento da perícia metrológica fotos, registro, cópia da fatura, quem assinou TOI, protocolo de entrega do TOI e comprovante de entrega do aviso, laudo do medidor etc.).
\item Solicita-se anulação do TOI referente a UC devido aos fatos narrados anteriormente.
\item Solicita-se a devolução dos valores do TOI em dobro, atualizados monetariamente pelo IPCA e juros de mora de 1\% ao mês calculados pro rata die, em conformidade com o estabelecido na Resolução Normativa n° 1.000/2021 da Aneel.
\item Efetue a devolução do valor total por meio de depósito bancário na conta corrente do Município, conforme o § 7º do art. 323 da Resolução Normativa nº 1000/2021 da ANEEL, e informe por escrito o resultado da análise, utilizando o endereço eletrônico: \EmailEmpresa{}.
\end{enumerate}
\end{enumerate}